\scanpage{21}
\mspage{19}

Mais peut-être ici faut-il un axiome plus subtil, p.ex. (si $\Sigma$ \inserted{\unsure{suspension|subséquents}} et $\Omega$ existent$^{*}$) que \deleted{les} \inserted{l'un des} morphismes d'adjonction $\Sigma\Omega\mathbin{\unsure{-|\longrightarrow}}\mathrm{id}$, $\mathrm{id}\longrightarrow\Omega\Sigma$ soit iso, ce qui implique que $\Omega$ resp. $\Sigma$ est pleinement fidèle. \deleted{Dans ce cas, on} En tous cas, l'attente

% slanted
\margin{\unsure{cf|v.} plus bas ce que ça signifie}

naturelle ici, c'est que $\mathrm{Ch}_{\infty}^{\Sigma}(a,b)$ soit non seulement dans $\mathrm{Hot}_{\bullet}$, mais même dans la catégorie \deleted{des complexes \unsure{objets}} \deleted{$\unsure{\mathrm D(k)}$} $\mathrm{D}^{\leq 0}(k\text{-}\mathrm{Mod})$, sous-catégorie pleine de la catégorie dérivée $\mathrm D^{-}(k\text{-}\mathrm{Mod})$ formée des $K$ tels que $H^i(K)=0$ si $i>0$. On peut aussi la noter $\mathrm D_{\geq 0}(k\text{-}\mathrm{Mod})$, l'interprétant en termes de complexes de ``chaînes'' plutôt que de cochaînes. \deleted{\unsure{Dans}}

\uline{Ici} les accouplements doivent correspondre à l'opération $*=\overset{\unsure{\mathbb L|L}}{\otimes}_{k}$ dans $\mathrm D_{\geq 0}(k\text{-}\mathrm{Mod})$, admettant $k$ comme objet unité, et

\[
 \unsure{\Pi|\pi}_0(K_\bullet)=H_0(K_\bullet).
\]

Donc on doit avoir \deleted{\unsure{un produit}}

\[
 \deleted{\unsure{\mathrm{RHom}}}\mathrm{Hom\,cat}_{\deleted{\unsure{\Sigma}}}(a,b)
 \overset{\mathbb L}{\otimes}_k\mathrm{Hom\,cat}(b,c)\longrightarrow\mathrm{Hom\,cat}(a,c),
\]

où cette fois $\mathrm{Hom\,cat}$ désigne l'objet de $\mathrm D_{\geq 0}(k\mathrm{Mod})$, au lieu d'un objet de $\mathrm{Hot}$ seulement, ou de $\mathrm{Hot}_{\bullet}$. \unsure{Alors utiliser|On utilise} les foncteurs

\[
 \mathrm D_{\geq 0}(k\,\mathrm{Mod})\longrightarrow\mathrm{Hot}_{\unsure{\bullet|*}}\longrightarrow\mathrm{Hot}
\]
% archive sheet number: 20
